\section{Conclusion}
\label{sec:conclusion}

We presented the first empirical comparison of multiprecision
(\burgerdybvig) and fixed-width (\schubfach) digit generation inside a
\rfcjcs/\ecma-constrained JSON canonicalization pipeline.  Our
conformance-gated benchmark methodology ensures that performance
measurements are only meaningful for correct implementations.

Key findings:
\begin{itemize}
\item \schubfach achieves 1.3--1.7$\times$ speedup on number-heavy
      workloads in API benchmarks (1.1--1.3$\times$ end-to-end CLI)
      due to zero heap allocation.
\item The advantage narrows to $<$5\% on string-dominant inputs.
\item Both implementations produce byte-identical output across x86-64
      and arm64.
\item Digit generation accounts for 15--40\% of total canonicalization
      cost depending on number density.
\end{itemize}

The benchmark infrastructure, 286{,}362 oracle vectors, and all raw data
are released as a reproducibility artifact.

\subsection*{Future Work}

Bare-metal arm64 performance benchmarks, additional language
implementations (Rust, C), multi-engine oracle validation (SpiderMonkey,
JavaScriptCore), and workload characterization from production JSON
corpora.
